% N-a-002, original Rado1964 presentation, including the transfinite section.
\chapter{Rado's Universal Graphs and Functions}
\label{sec:p03-rado-1964}

Universality depends on the class of objects and on the maps allowed between
them. Rado's 1964 paper constructs universal countable graphs, uniform and
mixed complexes, and finite-colour functions by coding finite requirements
into integer digits. It also gives a nonuniversality theorem for locally finite
graphs and a separate transfinite function theorem under specified cardinal
arithmetic assumptions \cite[pp.~331--340]{Rado1964}.

The constructions retain these different domains. An induced graph embedding
preserves edges and nonedges; a colour-function embedding acts on the arguments
while leaving colours fixed. The transfinite section allows repeated earlier
arguments and a final maximum, unlike the strictly increasing finite tuples
used in the digit constructions. The ambient set theory is ZFC. Additional
GCH assumptions occur only where stated.

\section{Graphs, weak embeddings and binary digits}
\label{sec:p03-rado64-graphs}

\begin{hbdefinition}{def:p03-rado64-graph}{The source graph and embedding conventions}
A graph is a pair \(\Gamma=(A,B)\), where \(B\subseteq[A]^2\);
\([A]^r\) denotes the set of \(r\)-element subsets of \(A\).
Thus graphs here are simple and undirected. A weak embedding is an injection
\(\varphi:A\to A^*\) preserving edges. Write \(xE_\Gamma y\) for
\(\{x,y\}\in B\). An embedding additionally reflects
edges:
\[
 \{x,y\}\in B\quad\Longleftrightarrow\quad
       \{\varphi(x),\varphi(y)\}\in B^*\qquad(x\ne y).
\]
This is an induced embedding. A graph universal for a class belongs to the
class and admits the indicated embedding of every member of that class.
\end{hbdefinition}

The two vertices of an edgeless graph weakly embed into a single edge,
but do not embed as an induced subgraph. Replacing the embedding notion
changes the question \cite[1.1]{Rado1964}.

\begin{hbdefinition}{def:p03-rado64-digits}{Digits and the binary graph}
Write \(\mathbb N_0=\{0,1,2,\ldots\}\). For an integer base \(k\geq2\),
let \(d_k(t,j)\in\{0,\ldots,k-1\}\) be the \(j\)-th digit of \(t\in\mathbb N_0\):
\[
 t=\sum_{j\geq0}d_k(t,j)k^j.
\]
Only finitely many digits are nonzero. Rado's graph \(R\) has vertex set
\(\mathbb N_0\) and, for \(x<y\),
\[
 xE_Ry\quad\Longleftrightarrow\quad d_2(y,x)=1.
\]
Adjacency is extended symmetrically, with no loops.
\end{hbdefinition}

A nonzero \(x\)-th binary digit of \(y\) implies \(y\geq2^x>x\).
This explains the orientation implicit in the original definition using
the two-element set \(\{x,y\}\).

\begin{hbtheorem}{thm:p03-rado64-graph-universal}{proved=annex:p03.rado64.graph-universal}{Rado's Theorem 1: an explicit induced universal graph}
Every graph with at most countably many vertices embeds into \(R\).
For a graph labelled by \(\mathbb N_0\), an increasing embedding is obtained
recursively. At stage \(\xi\), choose a nonnegative integer \(\psi(\xi)\) larger than every
previous \(\varphi(x)\), \(x<\xi\), and put
\[
 \varphi(\xi)=2^{\psi(\xi)}
       +\sum_{\substack{x<\xi\\xE_\Gamma\xi}}2^{\varphi(x)}.
\]
The unused high digit makes the map increasing without changing its required
adjacency digits.
\end{hbtheorem}

The recursion is effective relative to the input graph's adjacency information;
arbitrary countable graphs need not have computable adjacency. Universality
does not require surjectivity or imply uniqueness among all universal graphs.

\section{The extension property and the Fra\"iss\'e interface}
\label{sec:p03-rado64-extension}

\begin{hbtheorem}{thm:p03-rado64-extension}{proved=annex:p03.rado64.extension}{Finite adjacency patterns and the countable extension graph}
For disjoint finite \(U,V\subseteq R\), infinitely many vertices outside
\(U\cup V\) are adjacent to all of \(U\) and to none of \(V\).
Any two countable graphs with this property are isomorphic, and every
isomorphism between finite induced subgraphs of one extends to an automorphism.

In particular \(R\) is isomorphic to the graph in
\cref{thm:p03-fraisse54-graph}. Its extension axioms form the same complete
\(\aleph_0\)-categorical theory. These are the explicit modern graph
consequences, not additional assertions attributed to Rado's Theorem 1.
\end{hbtheorem}

\begin{hbtheorem}{thm:p03-rado64-amalgams}{proved=annex:p03.rado64.amalgams}{A graph amalgam and an amalgamation failure}
Finite simple graphs have the hereditary, joint-embedding and amalgamation
properties for induced embeddings. Their free amalgam identifies the common
induced graph and adds no edges between the two new sides.

Finite forests fail amalgamation for induced embeddings. Take a common graph
of two nonadjacent vertices \(a,b\), one extension \(a-x-b\) and another
\(a-y-z-b\), each an induced path. These have no common forest amalgam.
\end{hbtheorem}

A structural embedding is not by definition a partial elementary map.
The first-order graph conclusions use the separately supplied extension-axiom
argument. The deterministic binary graph is not a distribution on finite graphs.

\section{Why locally finite graphs have no universal member}
\label{sec:p03-rado64-locally-finite}

\begin{hbdefinition}{def:p03-rado64-local}{Local finiteness and finite-step reachability}
A graph is locally finite if each vertex has finitely many neighbours.
For \(m\geq1\), let \(k(\Gamma,x,m)\) count vertices reachable from \(x\)
by a walk of exactly \(m\) edges; vertices in a walk may repeat.
This count is finite in a locally finite graph.

Rado calls an acyclic graph a tree. The specific tree below is also connected,
so the distinction between tree and forest does not affect the construction.
\end{hbdefinition}

\begin{hbtheorem}{thm:p03-rado64-no-local-universal}{proved=annex:p03.rado64.no-local-universal}{Theorem 2, communicated by N. G. de Bruijn}
The class of countably infinite locally finite graphs has no weakly universal
member. For every countably infinite locally finite graph \(H\), there is
a countably infinite locally finite tree that does not weakly embed into \(H\).
\end{hbtheorem}

The candidate host must itself be locally finite. Rado's binary graph has
infinite vertex degrees, so there is no contradiction with Theorem 1.

\section{The binomial code for fixed-size sets}
\label{sec:p03-rado64-binomial}

\begin{hbdefinition}{def:p03-rado64-combinadic}{Rado's increasing-tuple code}
For an integer \(\ell\geq1\) and \(0\leq x_0<\cdots<x_{\ell-1}\), put
\[
 \tau(x_0,\ldots,x_{\ell-1})
       =\sum_{i=0}^{\ell-1}\binom{x_i}{i+1}.
\]
Use \(\binom ab=0\) for integers \(0\leq a<b\).
For the empty tuple put \(\tau(\varnothing)=0\).
Colexicographic order compares tuples at their largest coordinate of
disagreement.
\end{hbdefinition}

\begin{hbtheorem}{thm:p03-rado64-combinadic}{proved=annex:p03.rado64.combinadic}{The lemma of Section 2.2}
For each fixed \(\ell\geq1\), \(\tau\) is an order-preserving bijection from
increasing \(\ell\)-tuples of nonnegative integers in colexicographic order
onto \(\mathbb N_0\). Distinct \(\ell\)-element sets receive distinct digit
positions.
\end{hbtheorem}

This is the original fixed-dimension code. Mixed dimensions use a different
finite-subset code.

\section{Uniform and mixed complexes}
\label{sec:p03-rado64-complexes}

\begin{hbdefinition}{def:p03-rado64-complex}{The source's complex conventions}
An \(\ell\)-dimensional complex is \(C=(A,B)\) with
\(B\subseteq[A]^{\ell+1}\). An embedding is an injection preserving and reflecting
membership of every \((\ell+1)\)-element subset.
A mixed complex has
\[
 B\subseteq[A]^{<\omega}\setminus\{\varnothing\},
\]
allowing unbounded finite simplex sizes. No downward-closure condition is
imposed. The term is used in Rado's specified sense, not replaced by another
convention for abstract simplicial complexes.
\end{hbdefinition}

\begin{hbtheorem}{thm:p03-rado64-uniform-complex}{proved=annex:p03.rado64.uniform-complex}{Rado's Theorem 3}
For \(\ell\geq1\), define \(C_\ell^*\) on \(\mathbb N_0\) by
\[
 \{x_0<\cdots<x_\ell\}\in B_\ell^*
 \quad\Longleftrightarrow\quad
 d_2(x_\ell,\tau(x_0,\ldots,x_{\ell-1}))=1.
\]
It is universal for at most countable complexes of this dimension.
For \(\ell=0\), the analogous universal complex takes the odd vertices
as its singleton simplices.
\end{hbtheorem}

\begin{hbdefinition}{def:p03-rado64-mixed-code}{The source's mixed-length code}
For an increasing tuple \(0\leq x_0<\cdots<x_{\ell-1}\), with
\(\ell\geq0\), put
\[
 \rho(x_0,\ldots,x_{\ell-1})=\sum_{i<\ell}2^{x_i}.
\]
The empty tuple has code zero. Identifying a finite subset with its unique
increasing enumeration, binary expansion makes \(\rho\) a bijection from
all finite subsets of \(\mathbb N_0\) onto \(\mathbb N_0\).
The code \(\tau\) is injective only after the tuple length has been fixed;
\(\rho\) distinguishes lengths as well as the entries.
\end{hbdefinition}

\begin{hbtheorem}{thm:p03-rado64-mixed-complex}{proved=annex:p03.rado64.mixed-complex}{Rado's Theorem 4}
The mixed complex on \(\mathbb N_0\) given by
\[
 \{x_0<\cdots<x_\ell\}\in B_{\rm mix}^*
 \quad\Longleftrightarrow\quad
 d_2(x_\ell,\rho(x_0,\ldots,x_{\ell-1}))=1,\qquad\ell\geq0,
\]
is universal for at most countable mixed complexes.
The singleton case uses digit zero of its vertex.
\end{hbtheorem}

The source excludes the empty simplex. If membership of the empty set were
an additional freely varying truth value, a single host could not represent
both values: every embedding maps the empty set to itself.

\section{Universal finite-colour functions}
\label{sec:p03-rado64-finite-functions}

\begin{hbdefinition}{def:p03-rado64-colour-functions}{The classes \(K_{k\ell}\) and \(K_k\)}
For integers \(k\geq2\), \(\ell\geq0\), \(K_{k\ell}\) consists of functions
on increasing \((\ell+1)\)-tuples of \(\mathbb N_0\) with values in
\(\{0,\ldots,k-1\}\). The class \(K_k\) uses all nonempty finite increasing
tuples with the same colour set.
A universal member \(f^*\) represents every \(f\) through one injection
\(\varphi:\mathbb N_0\to\mathbb N_0\), preserving the argument domain, such that
\[
 f(x_0,\ldots,x_\ell)=
       f^*(\varphi(x_0),\ldots,\varphi(x_\ell)).
\]
The constructed maps are increasing. Arguments are mapped, while colours
remain fixed.
\end{hbdefinition}

\begin{hbtheorem}{thm:p03-rado64-finite-colours}{proved=annex:p03.rado64.finite-colours}{Rado's Theorem 5}
For \(\ell\geq1\), a universal member of \(K_{k\ell}\) is
\[
 f_{k\ell}^*(x_0,\ldots,x_\ell)
       =d_k(x_\ell,\tau(x_0,\ldots,x_{\ell-1})).
\]
For \(\ell=0\), use \(f_{k0}^*(x)=d_k(x,0)\).
This unary clause is included as a separately proved boundary extension;
the printed Theorem 5(a) assumes \(\ell\geq1\).
A universal member of \(K_k\) is
\[
 f_k^*(x_0,\ldots,x_\ell)=
       d_k(x_\ell,\rho(x_0,\ldots,x_{\ell-1})).
\]
All representing injections can be chosen increasing.
\end{hbtheorem}

The binary complexes are the \(k=2\) cases. Rado also interprets colours as
coefficients modulo \(k\); a chosen increasing vertex order then determines
the canonical order for oriented simplices, preserved by the embeddings.
No face-closure or chain-complex axiom is inferred.

These classes are not the class of all algebraic operations on all ordered
tuples. Repetitions, arbitrary order and infinitely many colours require
their own argument domains and proofs.

\section{The transfinite domain and index notation}
\label{sec:p03-rado64-transfinite-data}

\begin{hbconvention}{def:p03-rado64-transfinite-indices}{Initial ordinals and Rado's cofinality indices}
Write \(\omega_n\) for the initial ordinal of cardinality \(\aleph_n\).
Rado's index-valued notation \(\mathrm{cf}(n)\) is the ordinal satisfying
\[
 \operatorname{cf}(\aleph_n)=\aleph_{\mathrm{cf}(n)}.
\]
The left-hand \(\operatorname{cf}\) is ordinary cardinal cofinality;
\(\mathrm{cf}(n)\) on the right is its aleph index, not ordinary
cofinality of the ordinal \(n\). The source also uses \(n-1=r\) if
\(n=r+1\), and \(n-1=n\) if \(n\) has no immediate predecessor.
In this convention the displayed \(n-1\) is a predecessor notation,
including at zero and limit ordinals, rather than ordinal subtraction.
The ordinary cardinal and cofinality conventions are those of
\cref{entry:p06-source-30,def:p06-cofinality-regular-singular}.
\end{hbconvention}

\begin{hbdefinition}{def:p03-rado64-transfinite-domain}{The class \(F(l,m,S)\)}
Let \(l,m\) be ordinals and \(S\ne\varnothing\).
For every \(\lambda<l\), the argument sequences
\(\bar x=(x_i)_{i\leq\lambda}\) have length \(\lambda+1\) and satisfy
\[
 x_i\leq x_\lambda<m\quad(i<\lambda).
\]
Earlier coordinates need not increase or be distinct, and may equal the
final coordinate. \(F(l,m,S)\) is the set of functions on the disjoint union
of these domains, with values in \(S\).
Its universal \(f^*\) represents every \(f\) by one injection \(\varphi:m\to m\)
that preserves the domain and satisfies
\[
 f((x_i)_{i\leq\lambda})=f^*((\varphi(x_i))_{i\leq\lambda}).
\]
For \(l=0\) there are no argument tuples. The constructions below use
increasing injections on the stated initial-ordinal domains; no transfer
to arbitrary ordinal order types of the same cardinality is assumed.
\end{hbdefinition}

\begin{hbtheorem}{thm:p03-rado64-order-type-boundary}{proved=annex:p03.rado64.order-type-boundary}{The argument order type cannot be replaced by its cardinality}
The class \(F(2,\omega+1,\{0,1\})\) has no universal member,
although \(F(2,\omega,\{0,1\})\) has one by the countable case
of Theorem 6 below. The two argument domains have the same cardinality.
\end{hbtheorem}

\begin{hbtheorem}{thm:p03-rado64-restriction}{proved=annex:p03.rado64.restriction}{Reducing arities and the colour set}
If \(F(l,m,S)\) has a universal member, then so does
\(F(l_0,m,S_0)\) for \(l_0\leq l\) and
\(\varnothing\ne S_0\subseteq S\).
The reduction uses restriction and a colour-set retraction; the original
universal function need not already take its values in \(S_0\).
\end{hbtheorem}

\begin{hbdefinition}{def:p03-rado64-marked-template}{A template for a new maximum}
For a bound \(\eta\) in an ordered domain put
\(D_\eta=\{x:x<\eta\}\sqcup\{*\}\), with a new marker \(*\).
An \(l\)-template is a function
\[
 h:\bigsqcup_{\lambda<l}D_\eta^\lambda\longrightarrow S.
\]
Suppose each vertex \(z\) has been assigned one pair \((\eta_z,h_z)\).
For a tuple \((y_i)_{i\leq\lambda}\) ending at \(z=y_\lambda\), define
\[
 f^*((y_i)_{i\leq\lambda})=h_z((u_i)_{i<\lambda}),\qquad
 u_i=\begin{cases}y_i,&y_i<\eta_z,\\ *,&y_i\geq\eta_z.\end{cases}
\]
The tagged disjoint union retains the prefix length, even for an empty
prefix. This supplies a total function on the stated domain once the
assignments are fixed.
In the embedding construction, \(*\) denotes arguments equal to the newly
embedded maximum; all earlier images have already been placed below \(\eta\).
\end{hbdefinition}

\section{The three transfinite regimes}
\label{sec:p03-rado64-transfinite-theorem}

\begin{hbtheorem}{thm:p03-rado64-transfinite}{proved=annex:p03.rado64.transfinite}{Rado's Theorem 6 with the non-strict final-maximum domain}
Let \(l,n\) be ordinals and assume
\[
 2^{\aleph_\nu}=\aleph_{\nu+1}\qquad(\nu<n).
\]
The class \(F(l,\omega_n,S)\) has a universal member in each case:
\begin{enumerate}
\item \(n=n-1=\mathrm{cf}(n)\), \(l<\omega_n\), \(1\leq|S|\leq\aleph_n\);
\item \(n=r+1\), \(l<\omega_{\mathrm{cf}(r)}\), \(1\leq|S|\leq\aleph_n\);
\item \(n>\mathrm{cf}(n)\), \(l<\omega_n\), \(1\leq|S|<\aleph_n\).
\end{enumerate}
Its representing injections may be chosen increasing.
\end{hbtheorem}

In the first case the target is regular and either countable or a limit
cardinal. In the second it is a successor, with arity bound below its
predecessor's cofinality. In the third it is singular and the number of
colours has a strict upper bound. No GCH assumption entered the integer-digit
constructions.

The proof makes repeated occurrences of the final maximum explicit using
the marker in \cref{def:p03-rado64-marked-template}. Thus its requirements
are specified before the new image is chosen. Cardinal estimates and
regular/singular scheduling then establish the full stated non-strict domain.

Rado closes by asking how essential the transfinite argument restrictions
and successor-length tuples are, and for nontrivial nonuniversality cases.
These are the paper's historical questions, not assertions of present open
status or permission to remove the theorem's restrictions.

\begin{hbtheorem}{thm:p03-rado64-arity-obstruction}{proved=annex:p03.rado64.arity-obstruction}{Why all finite repeated-argument arities change the problem}
The class \(F(\omega,\omega,\{0,1\})\) has no universal member.
This is an additional explicit boundary example: it does not contradict
Theorem 5, whose arguments are strictly increasing and hence distinct.
\end{hbtheorem}

\section{Choosing the relevant universal object}
\label{sec:p03-rado64-reference}

The binary graph addresses induced embeddings of countable simple graphs.
The locally finite result changes the host class and fails even for weak
embeddings. Uniform complexes fix subset size; mixed complexes allow all
finite nonempty sizes. Finite-colour functions preserve colour values, while
the transfinite theorem uses a different sequence domain and additional
cardinal hypotheses.

For graphs, specify reflection of nonedges, the host class and countability.
For functions, specify arity, repetition, order, colours and the representation
map. The common word \emph{universal} cannot replace these distinctions.
